\pagelayout{margin}% Use a 1.5 column layout
\setchapterstyle{kao} % Choose the default chapter heading style
\chapter*{前言}
% \addcontentsline{toc}{chapter}{前言} % Add the preface to the table of contents as a chapter
\AddBackground

\marginnote[-3cm]{↑这张图片来自于\href{https://unsplash.com/photos/grayscale-photo-of-trees-covered-by-clouds-QSwk-0XxDBk}{Elias Maurer在Unsplash上的照片}}
\noindent 新年快乐！

此前很早就在网络上看到过KaoBook类的介绍，但是一直没有机会使用。但在实际编译和使用过程中，总是发现与中文适配的问题。不管是页面设计，还是minitoc的设计，虽然非常出彩，但实际使用过程中，发现稍加变化就会导致排版上的种种问题。直到最近看到在\LaTeX{} Studio上的更新，才发现薛浩先生已经将该模板进行了非常好的汉化工作并更新了\href{https://ask.latexstudio.net/ask/article/747.html}{Kaobook的汉化版}。有了这个前提，才让本文以及该模板有了可能，因此，我仅在此处将该模板进行了一些小的修改，以便适配我最近看到的一个非常好的排版效果。

\begin{itemize}
	\item \href{https://ask.latexstudio.net/ask/article/747.html}{薛浩在\LaTeX{} Studio上的更新}，其中包含了对Kaobook的汉化工作。
	\item \href{https://ask.latexstudio.net/addons/ask/go/index?url=https%3A%2F%2Fgithub.com%2Fxuehao%2FkaobookCJKsc}{薛浩在Github上的仓库地址}。
\end{itemize}


本模板是近期看到的一个非常好的排版效果，我看到后感觉这个论文当时一定就是用\LaTeX 的KOMA-Script类来完成的，因此，我在薛浩汉化版的基础上，再次进行了一些小的修改，以便适配中文以及复刻了这个排版效果，希望对大家后续使用有所帮助。

在制作本模板的时候，查阅了KaoBook类以及KOMA-Script的手册，不得不说，KOMA-Script是一个一场强大和精美的模板，但复杂程度和学习成本也是相当高的，官方的手册德议英的正文就有500余页。因此，这里也希望大家能够多多学习KOMA-Script的使用和交流。

\begin{flushright}
	\textit{Benjamin}\\
	\zhtoday
\end{flushright}
